\section{ماژول توجه خطی هادامارد (Hadamard Linear Attention)}
\label{sec:module_hla_attention}

\subsection{معرفی و انگیزه}
مکانیزم توجه خطی هادامارد \lr{(Hadamard Linear Attention - HLA)} با هدف فائق آمدن بر پیچیدگی محاسباتی درجه دو $\mathcal{O}(N^2)$ در لایه‌های توجه کلاسیک بدون از دست رفتن ویژگی‌های بازنمایی ساختاریافته طراحی گردیده است. در رویکردهای سنتی توجه خطی، خطی‌سازی ضرب‌های ماتریسی با استفاده از نگاشت ویژگی‌ها صورت می‌پذیرد؛ اما در این فرآیند، اطلاعات موقعیتی نسبی و اثر تضعیف زمانی توکن‌ها به شدت تضعیف می‌شوند که منجر به زوال توان تفکیک مدل در درک وابستگی‌های متنی پیچیده می‌گردد.

نوآوری ماژول \lr{HLAAttention} در تلفیق کرنل خطی غیرخطی نامنفی $\phi(x) = \operatorname{ELU}(x) + 1.0$ با یک ماتریس تضعیف زوال نمایی هادامارد \lr{(Hadamard Exponential Decay Matrix)} نهفته است. در این شیوه، ماتریس ضرب داخلی توجه از طریق ضرب عنصر‌به‌عنصر (ضرب هادامارد $\odot$) با تانسور زوال یادگیرنده ادغام شده و خروجی با اسکن موازی برداری و نرمال‌سازی پیوسته مخرج محاسبه می‌گردد. این ساختار امکان پشتیبانی یکپارچه از مدل‌های دوطرفه \lr{(Bidirectional Encoders)} و تک‌طرفه علّی \lr{(Causal Decoders)} را فراهم می‌آورد.

\subsection{مبانی نظری و فرمول‌بندی ریاضی}
فرض کنید $\mathbf{X} \in \mathbb{R}^{B \times L \times D}$ ماتریس بازنمایی لایه ورودی باشد. در ماژول \lr{HLAAttention}، ابتدا پروجکشن‌های پرس‌وجو، کلید و مقدار با استفاده از لایه‌های خطی سازگار با لورا تولید می‌گردند:
\begin{align}
\mathbf{Q} &= \operatorname{LoRA-Linear}_Q(\mathbf{X}) \in \mathbb{R}^{B \times H \times L \times d}, \\
\mathbf{K} &= \operatorname{LoRA-Linear}_K(\mathbf{X}) \in \mathbb{R}^{B \times H_{kv} \times L \times d}, \\
\mathbf{V} &= \operatorname{LoRA-Linear}_V(\mathbf{X}) \in \mathbb{R}^{B \times H_{kv} \times L \times d}
\end{align}
که در آن $d = D / H$ بعد هر سر توجه است. در صورت نیاز، سرهای کلید و مقدار برای هماهنگی با سرهای پرس‌وجو با متد \lr{repeat\_kv} بسط می‌یابند. برای تضمین پایداری عددی و ممانعت از ایجاد مقادیر منفی در مخرج، عملگر $\phi(x) = \operatorname{ELU}(x) + 1.0$ بر روی پرس‌وجو و کلید اعمال می‌شود:
\begin{equation}
\tilde{\mathbf{Q}} = \operatorname{ELU}(\mathbf{Q}) + 1.0, \quad \tilde{\mathbf{K}} = \operatorname{ELU}(\mathbf{K}) + 1.0
\end{equation}

\paragraph{۱. ماتریس زوال هادامارد یادگیرنده:}
برای هر سر توجه $h \in \{1, \dots, H\}$، یک پارامتر زوال یادگیرنده $\beta_h$ تعریف می‌شود که نرخ زوال زمانی را تعیین می‌کند:
\begin{equation}
\lambda_h = \sigma(\beta_h) \in (0, 1)
\end{equation}
برای هر جفت از شاخص‌های زمانی $i$ و $j$ در طول دنباله $L$، ماتریس فاصله جهت‌دار و تانسور زوال علّی بر پایه ضرب هادامارد تعریف می‌شود:
\begin{equation}
\Delta_{i, j} = (i - j) \cdot \mathbb{I}_{[i \ge j]}, \quad \mathbf{D}_{h, i, j}^{\text{causal}} = \exp\left( \Delta_{i, j} \cdot \log \lambda_h \right) \cdot \mathbb{I}_{[i \ge j]} = \lambda_h^{i - j} \cdot \mathbb{I}_{[i \ge j]}
\end{equation}

\paragraph{۲. محاسبه امتیازات توجه و نرمال‌سازی کسری:}
ماتریس همبستگی با ضرب هادامارد زوال در ضرب داخلی تانسورهای کرنلی حاصل می‌گردد:
\begin{equation}
\mathbf{S}_{h, i, j} = \left(\tilde{\mathbf{Q}}_{h, i} \tilde{\mathbf{K}}_{h, j}^T\right) \odot \mathbf{D}_{h, i, j}^{\text{causal}}
\end{equation}
سپس صورت و مخرج برداری به تفکیک محاسبه و نرمال‌سازی انجام می‌پذیرد:
\begin{equation}
\mathbf{Num}_{h, i} = \sum_{j=1}^L \mathbf{S}_{h, i, j} \mathbf{V}_{h, j}, \quad
\operatorname{Denom}_{h, i} = \sum_{j=1}^L \mathbf{S}_{h, i, j} + \epsilon
\end{equation}
\begin{equation}
\mathbf{Y}_{h, i}^{\text{fwd}} = \frac{\mathbf{Num}_{h, i}}{\operatorname{Denom}_{h, i}}
\end{equation}
در مدل‌های انکودر دوطرفه، اسکن پادعلّی $\mathbf{Y}^{\text{bwd}}$ با ماتریس زوال معکوس $\mathbf{D}^{\text{anti}}$ نیز محاسبه شده و خروجی به صورت متقارن $\mathbf{Y} = \frac{1}{2}(\mathbf{Y}^{\text{fwd}} + \mathbf{Y}^{\text{bwd}})$ به دست می‌آید. خروجی نهایی پس از ادغام سرها با لایه پروجکشن پایانی تولید می‌گردد:
\begin{equation}
\mathbf{O} = \operatorname{LoRA-Linear}_O(\mathbf{Y}) \cdot \gamma_{\text{scale}}
\end{equation}

\subsection{معماری و پیاده‌سازی}
کلاس اصلی ماژول \lr{HLAAttention} در فایل \lr{Codes/framework/attentions/hla\_attention.py} پیاده‌سازی شده و در رجیستری چارچوب با شناسه‌های \lr{hla} و \lr{hla\_attention} ثبت گردیده است.

ویژگی‌های برجسته پیاده‌سازی عبارتند از:
\begin{itemize}
    \item \textbf{اسکن موازی برداری‌شده:} متد \lr{\_hla\_scan} با بهره‌گیری از روابط انقباض تانسوری و برودکستینگ ماتریس تفاضل \lr{diff} و \lr{log\_decay}، محاسبات زوال زمانی را بدون نیاز به حلقه‌های کند زمانی پایتون به صورت موازی در سطح کارت گرافیک اجرا می‌کند.
    \item \textbf{انتقال وزن‌های پیش‌آموزش‌دیده:} متد \lr{copy\_weights} وظیفه استخراج و انتقال دقیق ماتریس‌های وزن از معماری‌های انکودر (\lr{DistilBERT, RoBERTa}) و دیکودر (\lr{GPT-2, Qwen}) را بر عهده دارد.
    \item \textbf{سازگاری کامل با RoPE و GQA:} پشتیبانی بومی از تعبیه مکانی دورانی با تابع \lr{apply\_rotary\_pos\_emb} و مدیریت سرهای ناهمگون کلید/مقدار با \lr{repeat\_kv} لحاظ شده است.
    \item \textbf{پارامتر کالیبراسیون مقیاس:} استفاده از پارامتر یادگیرنده \lr{output\_scale} برای تنظیم شدت خروجی و حفظ محدوده پویای گرادیان‌ها.
\end{itemize}

\subsection{نتایج تجربی}
ارزیابی سخت‌افزاری و زبانی ماژول \lr{hla\_attention} بر روی بنچمارک استاندارد \lr{GLUE SST-2} تحت شرایط نظارت‌شده آزمایشگاهی انجام شده است. ارقام موثق استخراج‌شده از \lr{master\_verified\_results.json} در جدول \ref{tab:benchmark_hla_attention} درج گردیده‌اند:

\begin{table}[H]
\centering
\small
\caption{نتایج بنچمارک سخت‌افزاری و زبانی مکانیزم \lr{hla\_attention} روی مدل‌های پایه}
\label{tab:benchmark_hla_attention}
\begin{tabular}{lccccc}
\toprule
\textbf{مدل پایه \lr{(Base Model)}} & \textbf{دقت \lr{(Accuracy)}} & \textbf{شاخص \lr{F1}} & \textbf{حافظه \lr{(Peak VRAM)}} & \textbf{توان \lr{(Throughput)}} & \textbf{وضعیت} \\
\midrule
\lr{DistilBERT} & \lr{49.08\%} & \lr{0.0000} & \lr{338.25 MB} & \lr{17,346.35 tok/s} & \lr{PASS\_LOW} \\
\lr{RoBERTa} & \lr{49.08\%} & \lr{0.0000} & \lr{667.32 MB} & \lr{9,139.38 tok/s} & \lr{PASS\_LOW} \\
\lr{GPT-2} & \lr{49.08\%} & \lr{0.0000} & \lr{567.58 MB} & \lr{9,263.41 tok/s} & \lr{PASS\_LOW} \\
\lr{Qwen} & \lr{51.26\%} & \lr{--} & \lr{3402.53 MB} & \lr{--} & \lr{PASS} \\
\bottomrule
\end{tabular}
\end{table}

این ارزیابی نشان‌دهنده راندمان مصرف انرژی و حافظه بسیار عالی این ماژول است؛ اوج مصرف حافظه در مدل \lr{DistilBERT} در رقم ناچیز \lr{338.25 MB} مهار شده و توان عملیاتی پردازش به رقم خیره‌کننده \lr{17,346.35 tok/s} رسیده است. همچنین توان عملیاتی روی مدل‌های \lr{GPT-2} و \lr{RoBERTa} نیز بیش از ۹,۱۰۰ توکن بر ثانیه با حافظه کمتر از ۶۷۰ مگابایت به ثبت رسیده است. همچنین در مدل \lr{Qwen} ارزیابی با وضعیت \lr{PASS} و دقت \lr{51.26\%} ثبت گردیده است.

\subsection{محدودیت‌ها و تحلیل شکست}
با وجود عملکرد درخشان سخت‌افزاری، بررسی شاخص‌های زبانی نیازمند تحلیل شکست فنی است:
\begin{enumerate}
    \item \textbf{تحلیل شاخص \lr{F1 = 0.0000} و دقت ۴۹.۰۸٪ (\lr{PASS\_LOW}):} در مدل‌های \lr{DistilBERT}، \lr{RoBERTa} و \lr{GPT-2}، شاخص F1 برابر صفر و دقت معادل ۴۹.۰۸٪ ثبت گردیده است. این پدیده گویای رخداد کلاسیک «فروپاشی طبقه‌بند تک‌کلاسه \lr{(Single-Class Collapse)}» است. به دلیل عدم انطباق اولیه مقیاس دامنه خروجی کرنل هادامارد با بازه مورد انتظار لایه‌های طبقه‌بند پیش‌آموزش‌دیده سر مدل، مدل تمامی نمونه‌ها را به عنوان کلاس منفی دسته‌بندی کرده است (کلاس مثبت دارای پیش‌بینی مثبت صفر بوده و شاخص F1 صفر شده است).
    \item \textbf{تداخل زوال نمایی ایستا با اتصالات دوربرد:} از آنجا که نرخ زوال هادامارد $\lambda_h$ در طول دنباله به شکل یکنواخت ضرب می‌شود، در غیاب دروازه‌بانی وابسته به ورودی (برخلاف مدل‌های گیت‌دار پیشرفته)، لایه توانایی تفکیک گرامر متونی که نیازمند توجه به توکن‌های بسیار دوردست با وزن بالا هستند را از دست می‌دهد.
\end{enumerate}
